\section{What logarithmic comparison must still control}\label{sec:comparisons}

The approximation theorems are compatible with either a real or a
nonreal local zero configuration. To obtain a favorable-value
estimate one needs information about the canonical field beyond
the error of its approximation. We now make that distinction
quantitative for several natural comparison recipes.

\subsection{The exact signed-Jensen decomposition}

Let $\varphi\ge0$ be a nonzero compact $C^{1,1}$ test in a disk.
Its Laplacian has integral zero. Moreover its positive and
negative parts both have positive mass: otherwise $\Delta\varphi$
would vanish, and a compact harmonic function would be zero.
Set
\begin{equation}\label{eq:signed-measures}
 \beta_\varphi=\int(\Delta\varphi)_+\dd A
             =\int(\Delta\varphi)_-\dd A>0,\qquad
 \dd\mu_\pm=\frac{(\Delta\varphi)_\pm}{\beta_\varphi}\dd A.
\end{equation}
For a nonzero holomorphic $f$ define
\begin{equation}\label{eq:jensen-gaps}
 A_\pm=\int|f|^2\dd\mu_\pm,\qquad
 J_\pm=\log A_\pm-\int\log|f|^2\dd\mu_\pm\ge0.
\end{equation}
All integrals are finite; $A_\pm>0$ because holomorphic zeros have
area measure zero. Jensen's inequality follows by regularization
and a limit if necessary.

\begin{proposition}[Signed-Jensen identity]\label{prop:signed-jensen}
The zero statistic satisfies
\begin{equation}\label{eq:signed-jensen}
 \nu_f(\varphi)=\frac{\beta_\varphi}{4\pi}
       \left[\log\frac{A_+}{A_-}+J_--J_+\right].
\end{equation}
In addition
\begin{equation}\label{eq:quadratic-positive}
 A_+-A_-=\frac4{\beta_\varphi}
                       \int\varphi|f'|^2\dd A\ge0,
\end{equation}
with strict inequality if $f$ is nonconstant.
\end{proposition}
\begin{proof}
Subtract the two equalities in \eqref{eq:jensen-gaps} and
apply \eqref{eq:weak-onepoint}. Integration by parts and
$\Delta|f|^2=4|f'|^2$ give \eqref{eq:quadratic-positive}.
If $f$ is nonconstant, $f'$ cannot vanish on the open set on
which $\varphi>0$.
\end{proof}

The positive sign of the quadratic difference does not control
the signed difference of the Jensen gaps. For a zero-free
nonconstant function, such as $f(z)=\ee^{tz}$ with real $t\ne0$,
the left side of \eqref{eq:signed-jensen} is zero. Thus
$J_--J_+=-\log(A_+/A_-)<0$: even a zero-free field requires
exact cancellation of its quadratic variation.

Nor is $J_--J_+$ always nonpositive. On the unit disk take
$\varphi(z)=(1-|z|^2)_+^2$ and $f(z)=z^2+1/4$.
Writing $u=|z|^2$, one has
$\Delta\varphi=8(2u-1)$, $\beta_\varphi=2\pi$, and the
angular mean of $|f|^2$ is $u^2+1/16$.
Integration over $u>1/2$ and $u<1/2$ respectively gives
\[
 A_+=\frac{17}{24}+\frac1{16}=\frac{37}{48},\qquad
 A_-=\frac1{24}+\frac1{16}=\frac5{48}.
\]
The zeros $\pm\ii/2$ give
$\nu_f(\varphi)=2(1-1/4)^2=9/8$.
Since $9/8>\frac12\log(37/5)$, the gap difference is positive
here. The inequality follows already from the cubic Taylor lower
bound for $\ee^{9/4}>37/5$. This is a control on the generic
argument, not a canonical theta counterexample.

The quadratic quantities themselves have a literal source formula.
For $f=F/M_X$, substitute \eqref{eq:source} and
\eqref{eq:canonical-F} into $A_\pm$:
\begin{align}
 A_\pm&=M_X^{-2}\iint_{(0,\infty)^2}
               \Phi(u)\Phi(v)\mathcal K_\pm(u,v)\dd u\dd v,\label{eq:quadratic-source}\\
 \mathcal K_\pm(u,v)&=
       \int |D(z)|^2\cos(zu)\overline{\cos(zv)}\dd\mu_\pm(z).\notag
\end{align}
Fubini follows from the complete source majorant on the compact
disk. Each kernel defines a nonnegative quadratic form, but this
does not compare the logarithmic gaps.

With $g=\widehat F/M_X$ and $\varepsilon=\delta/M_X$,
$|f|,|g|\le1$ imply $|A_\pm-\widehat A_\pm|\le2\varepsilon$.
If $\widehat A_\pm>2\varepsilon$, the corresponding log error is at
most $2\varepsilon/(\widehat A_\pm-2\varepsilon)$.
Also \eqref{eq:EX} gives
\[
 \left|\int\log|f|^2\dd\mu_\pm-
               \int\log|g|^2\dd\mu_\pm\right|
 \le\frac{2\norm{\Delta\varphi}_\infty}{\beta_\varphi}E_X.
\]
These estimates acquire the separate quantities. They do not
bound their combination favorably.

\subsection{A sharp local distance from harmonic functions}

If a comparison field is zero-free, its logarithm is harmonic.
The following estimate measures the unavoidable absolute
logarithmic cost of using such a baseline near a zero.

\begin{theorem}[Sharp harmonic-distance bound]\label{thm:harmonic-distance}
Let $F$ be nonzero and holomorphic on $D(a,\ell)$, with a zero
of multiplicity $m\ge1$ at $a$. For every real harmonic function
$v$ on that disk,
\begin{equation}\label{eq:harmonic-distance}
 \int_{D(a,\ell)}|\log|F|-v|\dd A
 \ge\frac{\pi\log2}{2}\,m\ell^2 .
\end{equation}
The constant is sharp.
\end{theorem}
\begin{proof}
For $0<\varrho<\ell$, let $t=|z-a|$ and define
\[
 \Psi_\varrho(t)=
 \begin{cases}
 (\varrho^2\log2-t^2)/4,
                      &0\le t\le\varrho/\sqrt2,\\[2pt]
 (t^2-\varrho^2)/4-(\varrho^2/2)\log(t/\varrho),
                      &\varrho/\sqrt2\le t\le\varrho,\\[2pt]
 0,                   &t\ge\varrho.
 \end{cases}
\]
The value and first derivative agree at both junctions.
This is a nonnegative compact $C^{1,1}$ function, decreasing to
zero at $\varrho$, with
\[
 \Delta\Psi_\varrho=
 \begin{cases}-1,&t<\varrho/\sqrt2,\\
                1,&\varrho/\sqrt2<t<\varrho,\\
                0,&t>\varrho,\end{cases}
 \qquad
 \Psi_\varrho(0)=\frac{\varrho^2\log2}{4}.
\]
In particular $\norm{\Delta\Psi_\varrho}_\infty=1$, and there
are no boundary delta terms. The weak zero formula and harmonicity
give
\[
 \int(\log|F|-v)\Delta\Psi_\varrho\dd A
 =2\pi\sum_{F(z)=0}m_z\Psi_\varrho(z)
 \ge2\pi m\Psi_\varrho(0).
\]
Absolute values bound the left side by the $L^1$ distance.
Letting $\varrho\uparrow\ell$ proves
\eqref{eq:harmonic-distance}, also when that distance is infinite.

For sharpness take $F(z)=(z-a)^m$ and the constant harmonic
function $v=m\log(\ell/\sqrt2)$. With $u=|z-a|^2/\ell^2$,
the distance equals
\[
 \frac{\pi m\ell^2}{2}\int_0^1|\log(2u)|\dd u
 =\frac{\pi m\ell^2\log2}{2}.
\]
This matches the lower bound.
\end{proof}

Consider the precise recipe in which one estimates
$S_\varphi(F)=\nu_F(\varphi)$ by subtracting a harmonic $v$ and
paying the full-disk absolute allowance
\[
 B\ge\norm{\log|F|-v}_{L^1(D(c,r))},\qquad
 |S_\varphi(F)|\le K_\varphi B.
\]
If $X$ itself is a real zero of multiplicity $m$, the disk
$D(X,1/6)$ lies in $D(c,r)$, because $r-4=1/6$.
Theorem~\ref{thm:harmonic-distance} forces
\begin{equation}\label{eq:baseline-floor}
 B\ge\frac{\pi m\log2}{72}.
\end{equation}
For the scalar tests on their plateau, just above the real axis,
\[
 \Delta\varphi_Q=4,\qquad
 \Delta\varphi_\eta=
        \frac{4\eta(1-\eta y^2)}{(1+\eta y^2)^2}
                  \longrightarrow4\eta .
\]
Thus this recipe's allowances satisfy
\begin{equation}\label{eq:allowance-floors}
 K_QB\ge\frac{m\log2}{36},\qquad
 K_\eta B\ge\frac{\eta m\log2}{36}.
\end{equation}
They cannot pass the strict scalar consumers for
$d^2\le m\log2/72$, since
$2\log(1+\eta d^2)\le2\eta d^2$.
These are floors on the \emph{allowance}, not lower bounds on
the true transverse mass, which may be zero.
The excluded class is the harmonic-baseline comparison with this
full-disk absolute payment in a real-zero-centered window.
No conclusion follows about all weighted, signed, or real-divisor
comparisons.

\subsection{Keeping real-axis vanishing by regularization}

Quadratic vanishing of $\varphi$ does not make
$|\Delta\varphi|$ vanish at the axis. In particular the immediate
log-error bound is
$(2\pi)^{-1}\int|\Delta\varphi|\,|e|\dd A$,
not a $y^2$-weighted log norm. On a plateau,
$\Delta\varphi_Q=4$ for small positive $y$.
If arbitrary nonnegative residuals are concentrated in
$0<y<\epsilon$, a bound by
$C\int y^\beta |e|\dd A$, $\beta>0$, would require
$C\gtrsim\epsilon^{-\beta}$.
This only rules out the unqualified norm substitution; those
residuals are not asserted to be holomorphic log differences.

A legitimate way to retain the original test is to differentiate
a regularized logarithm. For $0<\tau\le1$ set
\[
 u_\tau(f)=\frac12\log(|f|^2+\tau^2),\qquad
 T_{\varphi,\tau}(f)=\frac1\pi
       \int\varphi\,\frac{\tau^2|f'|^2}
                            {(|f|^2+\tau^2)^2}\dd A .
\]
Holomorphy gives
$\Delta u_\tau=2\tau^2|f'|^2/(|f|^2+\tau^2)^2$, hence
$T_{\varphi,\tau}=(2\pi)^{-1}\int u_\tau\Delta\varphi\dd A$.
For $\varphi\ge0$ this integral is nonnegative. As $\tau\downarrow0$,
$u_\tau\downarrow\log|f|$ in local $L^1$, by domination with
$u_1-\log|f|$. Therefore
\begin{equation}\label{eq:regularized-limit}
 T_{\varphi,\tau}(f)\longrightarrow\nu_f(\varphi).
\end{equation}
Real multiple zeros receive their zero test weight in this limit;
nonreal zeros retain their positive weight.

The sublevel estimate \eqref{eq:sublevel} supplies a quantitative
bias when $n\ge2$. Write $L=L(A,n)$. Layer cake gives
\begin{align}
 \int(u_\tau(f)-\log|f|)\dd A
 &\le L\tau^{2/n}
       \int_0^\infty(\ee^{2v}-1)^{-1/n}\dd v\notag\\
 &\le 2nL\tau^{2/n}.\label{eq:regularization-bias}
\end{align}
Indeed $u_\tau-\log|f|>v$ means
$|f|<\tau(\ee^{2v}-1)^{-1/2}$.
On $0<v<1$, use $\ee^{2v}-1\ge2v$ to bound the last integral
by two; on $v\ge1$, use $\ee^{2v}-1\ge\ee^v$ to bound the
tail by $n$. Then $n+2\le2n$.
The canonical $n_X$ is at least two. Consequently
\[
 |T_{\varphi,\tau}(g)-\nu_g(\varphi)|
 \le B_{\varphi,\tau}
 :=2K_\varphi n_XL(A_X,n_X)\tau^{2/n_X}.
\]
This is a smoothing bias for one logarithm; its exponent
$2/n_X$ does not replace the source-transfer exponent
$2/(n_X+2)$.

The complete possible consumer would have to bound
\begin{equation}\label{eq:regularized-consumer}
 T_{\varphi,\tau}(g)+B_{\varphi,\tau}
       +K_\varphi E_X+E_{\rm eval}
       <\hbox{the declared defect margin}.
\end{equation}
The first term requires joint information about small values and
derivatives. The generic Cauchy bound from $R-r=1/3$ is
$|g'|\le3$, and yields only
$T_{\varphi,\tau}(g)\le9(\pi\tau^2)^{-1}\int\varphi\dd A$.
This allowance grows as the smoothing bias shrinks. The present
complete-source upper bounds therefore do not close this recipe.
This conclusion is scoped to those available inputs, not to
every signed or weighted logarithmic method.

\subsection{Removing a real divisor and accounting for what remains}

Let $P_R(z)=\prod_\ell(z-a_\ell)^{m_\ell}$ with real $a_\ell$.
For either scalar test, $\nu_{P_R}(\varphi)=0$.
For any harmonic $v$ one consequently has the identity
\begin{equation}\label{eq:real-divisor}
 \nu_F(\varphi)=\frac1{2\pi}\int\Delta\varphi
     \bigl(\log|F|-\log|P_R|-v\bigr)\dd A.
\end{equation}
It remains valid when $P_R$ is not a divisor of $F$: the
corresponding real poles have zero test weight.
To conclude that $F/P_R$ is holomorphic, however, every selected
root and its multiplicity must be justified independently.
Even removal of all real roots in a compact domain leaves any
nonreal roots in the quotient.

This representation does improve on the harmonic class for a
purely real divisor. If $F=(z-a)^m\ee^q$, with real $a$ and
holomorphic $q$, take $P_R=(z-a)^m$ and $v=\Re q$; the residual
is exactly zero. The harmonic allowance floor therefore does
not extend to this class.
Two costs remain. A field error cannot be divided by $P_R$
at its roots without shared vanishing or controlled deflation:
$F=z-a$, $\widehat F=z-a+\epsilon$ is an immediate example.
Also, the finite theta source has no acquired real-rootedness
theorem, so it cannot simply be declared an all-real comparator.

A concrete independent alternative would be a matched integer
count. Let a Jordan domain $\Omega$ contain the whole test
support, have closure within $D(c,r)$, and satisfy
$F\ne0$ on its boundary. Suppose rigorous integers give
\[
 N_\Omega=\frac1{2\pi\ii}\int_{\partial\Omega}\frac{F'}F\dd z
          \le U,\qquad
 \hbox{certified real multiplicity in }\Omega\ge M_R.
\]
The argument principle makes $N_\Omega$ the total multiplicity.
Thus $R_\Omega=U-M_R\ge0$ bounds all unaccounted zeros,
including every relevant nonreal zero. The canonical strip
then gives the valid coarse estimates
\begin{equation}\label{eq:count-only}
 Q\le2R_\Omega,\qquad
 \log Z(\eta)\le2\log(1+\eta)R_\Omega .
\end{equation}
If $R_\Omega\ge1$, these particular allowances are at least
the one-pair margin for every $0<d\le1$. The count-only recipe
therefore needs $U=M_R$ to pass the strict scalar consumer.
In that case all zeros in the domain are accounted for as real,
and the statistic is zero directly. Matching counts, a positive
boundary modulus, and real-axis sign or multiplicity certificates
have not been supplied here. Cauchy's estimate gives derivative
approximation error at most $3\delta$ on $D(c,r)$; it gives
no lower bound for $|F|$ on a proposed contour.

Finally, a common real divisor and a safe anchor cannot by
themselves force reality at fixed positive approximation tolerance.
The local controls
\begin{equation}\label{eq:real-divisor-control}
 f_\pm(z)=\frac{(z-X)^2\pm d^2}{5(-16\pm d^2)},
 \qquad 0<d\le1,
\end{equation}
have $f_\pm(c)=1/5$ and
$\norm{f_\pm}_\infty\le293/300<1$ on $D(c,R)$.
Direct subtraction yields
\[
 f_+-f_-=
 -\frac{2d^2((z-X)^2+16)}{5(256-d^4)},\qquad
 \norm{f_+-f_-}_\infty\le\frac{353}{2550}d^2.
\]
Their transverse masses are $2d^2$ and zero. Multiplying both
by a fixed real polynomial divided by its disk supremum
preserves a common known real divisor, an upper bound one,
and a common positive anchor modulus; all added real zeros
remain invisible to the test. These controls lack the canonical
theta source identity. They show why fixed generic local
information is insufficient, without excluding a theta-specific
estimate, scale-dependent precision, or a valid finite count
certificate.
